% =====================================================================
%  CHAPTER 1 -- INTRODUCTION
%
%  HEADING CONVENTION (see style/fypstyle.sty section 6):
%    \chapter{...}     and  \section{...}     are typed in CAPITALS
%    \subsection{...}  is typed in Sentence case
%
%  WHAT IS ALREADY WRITTEN HERE:
%  The background, problem statement, objectives and scope below are
%  CARRIED OVER FROM YOUR FYP PROPOSAL (sections 1, 1.1, 1.2, 3.1, 3.2,
%  4.2 and 6). They are your own words, reformatted -- not invented
%  content. Each block is marked so you can see what came from where.
%  Look for "% EXPAND:" comments showing where a final report needs more
%  than a proposal did.
%
%  This is also the first page of the Arabic page numbering (page 1).
% =====================================================================
\chapter{INTRODUCTION}
\label{ch:introduction}

% ---------------------------------------------------------------------
%  CHAPTER OPENING -- background and context.
%  CARRIED OVER FROM PROPOSAL SECTION 1 (Project Summary/Introduction).
% ---------------------------------------------------------------------

The architecture and construction industry, a cornerstone of urban
development, is currently hindered by significant fragmentation and
inefficiency. A persistent disconnect exists among key stakeholders:
clients with a vision, architects with the design expertise, and
construction companies with the execution capabilities. This discovery
gap forces clients to navigate a complex, manual landscape where
articulating their needs in technical terms is difficult, and finding
reliable professionals is often a gamble based on unverified referrals.

Simultaneously, talented architects and contractors struggle to showcase
their specialised portfolios to the right audience. They rely on
inefficient general marketplaces that capture neither the nuance of
design style nor the security required for high-stakes construction
agreements. The result is a pre-contract phase that is slow for every
party involved and trustworthy for none of them.

\fypProjectTitle{} addresses these systemic failures by introducing an
AI-augmented, three-sided marketplace that unifies the pre-construction
pipeline. By applying natural language processing to decode client
intent and semantic matching to align design sensibilities, the platform
replaces manual searching with intelligent recommendation. It fosters an
ecosystem in which trust is established through verified portfolios and
digital contract formalisation, so that every connection rests on
genuine compatibility and transparency. This simplifies the client's
journey from concept to contract while supplying professionals with a
stream of high-quality leads.

% EXPAND: a final report's introduction is normally one to two pages.
% Add a paragraph quantifying the problem -- the size of the AEC sector
% in Pakistan, or the typical duration and cost of the pre-contract
% phase. Cite a source for any figure you quote.

% ---------------------------------------------------------------------
\section{PROBLEM STATEMENT}
\label{sec:problem-statement}
% ---------------------------------------------------------------------
%  CARRIED OVER FROM PROPOSAL SECTION 1.1, with the justifying citations
%  added from Chapter 2 as the rubric requires.
%
%  RUBRIC NOTE (R3): the top band requires the problem statement to be
%  justified "with reference to the literature review". The citations
%  below do that. Do not remove them.
% ---------------------------------------------------------------------

The architecture and construction industry lacks a unified digital
ecosystem, making it difficult for clients to find professional
architects and verified contractors. While digital multi-sided platforms
have transformed comparable industries, the construction sector has no
equivalent ecosystem for stakeholder collaboration
\cite{hossain2018platform}.

Existing platforms fail to bridge the pre-contract gap in two specific
ways. First, the tools that apply natural language processing in this
domain operate on technical construction documents and cannot decode the
non-technical language in which a client actually expresses a design
intention \cite{lee2017extraction, wu2022nlp}. Second, the agreement
procedures that follow a match remain manual and opaque, which is the
mechanism by which fraud risk enters the process
\cite{li2019blockchain}.

There is therefore a need for an intelligent solution that enables
seamless collaboration, secure agreements, and intent-based design
discovery within a single platform.

% EXPAND: add one paragraph on the consequence of leaving this
% unsolved -- what clients, architects and contractors each lose today
% in time, money or risk. Quantify it if you can find a source.

% ---------------------------------------------------------------------
\section{OBJECTIVES}
\label{sec:objectives}
% ---------------------------------------------------------------------
%  CARRIED OVER FROM PROPOSAL SECTIONS 3.1 (Main Goal) and 3.2
%  (Specific Objectives).
% ---------------------------------------------------------------------

The main goal of this project is to develop an intelligent web platform
that connects architects, construction companies and the public for
efficient design discovery, communication and project collaboration. The
specific objectives are as follows:

% NOTE: the braces around a bracketed placeholder are REQUIRED -- without
% them LaTeX would read "[...]" as a custom label for \item.
\begin{enumerate}[label=\arabic*., leftmargin=*, itemsep=2pt]
  \item To allow architects to create and manage design portfolios with
        semantic tagging.
  \item To enable users to browse, filter and purchase designs, or
        contact architects directly, through natural language queries.
  \item To provide construction companies a channel through which to
        connect with potential clients using verified profiles.
  \item To integrate AI-based semantic search using S-BERT.
        % DISCREPANCY: the proposal's wording was "AI-based search and
        % voice query features using WhisperASR and S-BERT". Voice query
        % is not implemented. An objective may legitimately state intent
        % that was later descoped -- but if you keep the voice half here,
        % Chapter 6 must report honestly that it was not met. Cleaner to
        % state the objective as delivered, as above, and list voice
        % input under future work.
  \item To ensure secure authentication, data management and smooth
        communication between all stakeholders.
  \item To provide a real-time communication hub enabling secure
        interaction between clients, architects and contractors.
\end{enumerate}

% EXPAND: each objective above is functional. The rubric and your
% Chapter 6 results will both be stronger if two or three of these carry
% a MEASURABLE target -- for example "semantic search returning relevant
% results with a top-5 precision of at least X" or "voice transcription
% with a word error rate below Y". Add the targets you intend to
% evaluate against, because Chapter 6 has to report results for them.

% ---------------------------------------------------------------------
\section{SCOPE}
\label{sec:scope}
% ---------------------------------------------------------------------
%  ASSEMBLED FROM PROPOSAL SECTIONS 4.2 (Major Modules), 6.1 (Final
%  Product Includes), 6.2 (Platforms Supported) and 6.3 (APIs).
% ---------------------------------------------------------------------

This project delivers a responsive, browser-based web application
comprising five modules. The \textbf{User Module} handles registration,
authentication, browsing of architectural designs, and initiating
contact with architects. The \textbf{Architect Module} supports design
upload, portfolio management with semantic tagging, and direct
communication with clients. The \textbf{Construction Module} enables
construction companies to list services and connect with clients through
verified credentials. The \textbf{Admin Module} provides user
verification, data moderation and analytics. The \textbf{AI Module}
performs semantic matching of project briefs against architect
portfolios using Sentence-BERT embeddings \cite{reimers2019sbert}.

The delivered product includes interactive dashboards for each of the
three user types, a real-time chat feature with secure messaging,
semantic search over architect portfolios, digital agreements with
signing, and an administrative control panel with identity verification
workflows. The platform integrates the Stripe API for payment
processing and payouts.

% DISCREPANCY -- see the full note in Section~\ref{subsec:major-modules}
% of Chapter 3. The proposal also lists voice query input via WhisperASR
% and keyword extraction via SpaCy as part of the AI module. Neither is
% present in the repository, so neither is claimed here. If you
% implement them before submission, restore them to this paragraph and
% re-add the \cite{radford2022whisper} citation. If you do not, say so
% explicitly in the exclusions paragraph below and list them as future
% work in Section~\ref{sec:future-work}.

% TODO: WRITE THE EXCLUSIONS PARAGRAPH. Your proposal does not state
% what is out of scope, and this is the one part of Section 1.3 that
% cannot be carried over. Examiners actively look for it, and an honest
% exclusions paragraph reads as engineering maturity rather than as a
% weakness. Likely candidates, based on your proposal:
The following are outside the scope of this project. Native mobile
applications are excluded; the platform is delivered as a responsive
browser-based application. Construction project management after
contract signing is excluded, since that phase is already well served by
existing tools and this project targets the pre-contract gap
specifically. Handling of 3D or BIM model files, multi-language
support, and any on-site or IoT integration are also excluded.

% TODO: DECIDE AND THEN STATE IT. If voice query input is not
% implemented by submission, add a sentence here saying so -- for
% example: "Voice query input, proposed during the planning phase, was
% descoped during implementation and is discussed as future work in
% Section~\ref{sec:future-work}." Declaring a descoped feature costs you
% nothing; being caught with an undeclared one costs a great deal.
[VOICE INPUT: keep or delete this line once you have decided -- see the
discrepancy note in Section~\ref{subsec:major-modules}.]

% ---------------------------------------------------------------------
\section{REPORT ORGANIZATION}
\label{sec:report-organization}
% ---------------------------------------------------------------------
%  Boilerplate -- adjust only if you rename or merge chapters.
% ---------------------------------------------------------------------

The remainder of this report is organized as follows.

\textbf{Chapter~\ref{ch:literature-review}} reviews the existing
literature and related platforms, and identifies the research gap that
this project addresses.

\textbf{Chapter~\ref{ch:methodology}} presents the methodology,
describing the development approach, the requirements analysis, and the
tools and techniques used.

\textbf{Chapter~\ref{ch:system-design}} describes the system
architecture and the detailed system design, including the UML models
and the entity relationship model of the system.

\textbf{Chapter~\ref{ch:implementation-testing}} covers the
implementation of the system and the testing performed against it,
including the test cases and their outcomes.

\textbf{Chapter~\ref{ch:results}} presents and discusses the results
obtained from the evaluation of the system.

\textbf{Chapter~\ref{ch:conclusion}} concludes the report and outlines
directions for future work.
